\subsection{Sockets TCP/UDP, RPC, gRPC, Protocol Buffers y RMI}
\label{sec:nixon-comunicacion}

\subsubsection{Sockets y mensajes de aplicación}
Un socket permite que dos programas intercambien datos mediante una dirección y
un puerto. En esta práctica los nodos se comunican por localhost, pero cada uno
mantiene su propio reloj lógico. Un nodo puede atender solicitudes como servidor
y enviar otras como cliente; esos papeles dependen del intercambio.

TCP establece una conexión y entrega a la aplicación un flujo de bytes fiable y
ordenado. Sin embargo, el flujo no conserva los límites de cada mensaje. Tampoco
un reconocimiento del transporte demuestra que la aplicación receptora haya
procesado una solicitud~\cite{belinda_eddy2022_tcp}. Por eso el equipo utiliza un
acuse de aplicación con el identificador del mensaje y el reloj del receptor.

En Python una llamada a \texttt{recv} puede devolver solamente parte de los bytes
esperados. El programa debe acumularlos hasta completar su mensaje, o detectar
que la conexión se cerró~\cite{nixon_python_sockets}. El archivo
\path{src/sockets/protocol.py} resuelve este problema anteponiendo al JSON una
cabecera de cuatro bytes que indica la longitud del cuerpo. Primero lee esa
cabecera y después el cuerpo completo. La codificación UTF-8 permite conservar
el contenido del texto y medir su tamaño en bytes.

UDP trabaja con datagramas y no exige establecer previamente una conexión.
La aplicación debe estar preparada para datagramas perdidos, repetidos o fuera
de orden~\cite{nixon_microsoft_udp}. Si necesitara confirmar la recepción tendría
que definir sus propios acuses y controles. UDP sirve aquí para comparar modelos;
los mecanismos implementados y medidos son TCP con JSON y gRPC con Protobuf.

\subsubsection{RPC y gRPC}
Una llamada a procedimiento remoto, o RPC, presenta una operación de otro proceso
mediante una interfaz de llamada. Un representante local, denominado \emph{stub},
prepara los argumentos para enviarlos; el servidor ejecuta la operación y devuelve
el resultado. La apariencia de una función local facilita el uso, aunque sigue
existiendo comunicación entre espacios de memoria distintos~\cite{nixon_oracle_rmi}.

gRPC concreta este modelo mediante un contrato de servicio y código generado para
cliente y servidor. Permite llamadas unarias y distintas formas de transmisión
de secuencias de mensajes. El servicio de la práctica usa una llamada unaria:
\texttt{EnviarMensaje} recibe un \texttt{Mensaje} y devuelve un
\texttt{Acuse}~\cite{nixon_grpc_conceptos}. El cliente reutiliza el canal al mismo
destino; no necesita construir manualmente las tramas JSON de TCP.

La comparación experimental corresponde a dos formas de programar la comunicación:
un protocolo propio sobre TCP y un servicio gRPC sobre HTTP/2 y TCP. Por tanto,
gRPC también utiliza TCP en esta implementación. Zhu y colaboradores analizan
cómo las capas de serialización, RPC, HTTP y transporte añaden trabajo a una
llamada remota~\cite{nixon_zhu2025_rpc}. Ese análisis ayuda a entender por qué un
contrato más cómodo para el programador no garantiza un RTT menor.

Las llamadas remotas necesitan un plazo de espera y tratamiento de errores.
gRPC proporciona estados de terminación y plazos por llamada~\cite{nixon_grpc_conceptos}.
En nuestro cliente, un plazo agotado se registra como \texttt{timeout}; un acuse
con identificador incorrecto se considera un error. No se reintenta automáticamente
una solicitud, porque perder la respuesta no demuestra que el receptor no la haya
procesado. El identificador permite relacionar los registros, pero por sí solo
no elimina duplicados si en el futuro se añadieran reintentos.

\subsubsection{Protocol Buffers y contrato compartido}
Protocol Buffers describe mensajes mediante campos con tipos y números estables.
Su formato binario utiliza esos números para identificar los campos; los textos
se codifican en UTF-8. El número asignado a un campo identifica su significado en
el formato y no representa un valor de ejemplo~\cite{nixon_protobuf_encoding}.

El contrato \path{proto/messaging.proto} declara \texttt{message\_id},
\texttt{sender} y \texttt{message} como texto, y \texttt{timestamp} como
\texttt{int64}. Conserva así los cuatro campos acordados con TCP. El código
generado conoce esos tipos, pero el servidor todavía debe comprobar reglas del
experimento, como que el identificador no esté vacío y que el emisor sea un nodo
válido. El acuse devuelve el mismo identificador, el nombre del receptor y su
reloj de envío para actualizar el reloj del cliente.

Una carga útil de 100 bytes significa que el contenido de \texttt{message} ocupa
100 bytes en UTF-8. El mensaje serializado también incluye los otros campos y la
comunicación añade sus propias cabeceras. Por ello, igualar la carga útil entre
los mecanismos no significa igualar todos los bytes transmitidos.

Niswar y colaboradores compararon REST, GraphQL y gRPC en microservicios con
consultas de datos planos y anidados. En sus pruebas gRPC obtuvo menores tiempos
de respuesta~\cite{nixon_niswar2024_microservices}. Ese estudio no compara el
protocolo TCP/JSON de esta práctica: cambian las operaciones y el entorno. Sus
resultados sirven como antecedente, mientras que la elección para nuestro caso
se apoya en los CSV obtenidos por el equipo.

\subsubsection{RMI y comparación conceptual}
Java RMI permite invocar métodos de objetos alojados en otras máquinas virtuales
Java mediante un representante local. Su modelo conserva conceptos propios de
los objetos, incluidas las referencias a objetos remotos~\cite{nixon_oracle_rmi}.
Comparte con RPC la idea de ocultar parte del intercambio de mensajes, pero su
abstracción se centra en objetos del entorno Java. En esta práctica se implementa
un servicio gRPC en Python; RMI se estudia como otra forma de comunicación remota.

La tabla~\ref{tab:nixon-modelos} resume las responsabilidades de cada modelo.
Las filas de TCP y gRPC describen las decisiones de este repositorio; UDP y RMI
corresponden a sus características conceptuales documentadas anteriormente.

\begin{table}[htbp]
\centering
\small
\begin{tabular}{@{}p{2.1cm}p{3.5cm}p{4cm}p{4.7cm}@{}}
\toprule
Modelo & Unidad para el programa & Representación & Trabajo de la aplicación \\
\midrule
Socket TCP & Flujo de bytes conectado & JSON UTF-8 con longitud & Delimitar y leer mensajes, validar campos y acuses, controlar errores. \\
\addlinespace
Socket UDP & Datagrama & Bytes elegidos por la aplicación & Prever pérdidas, duplicados y desorden según sus necesidades. \\
\addlinespace
gRPC & Operación de un servicio & Mensajes Protobuf en este proyecto & Definir contrato, implementar métodos y atender estados y plazos. \\
\addlinespace
Java RMI & Método de un objeto remoto & Objetos y referencias del entorno Java & Definir interfaces remotas, implementar objetos y atender fallos remotos. \\
\bottomrule
\end{tabular}
\caption{Comparación conceptual de los modelos de comunicación.}
\label{tab:nixon-modelos}
\end{table}

\subsubsection{Secuencia de un intercambio}
Las figuras~\ref{fig:nixon-secuencia-tcp} y~\ref{fig:nixon-secuencia-grpc}
representan el recorrido de una solicitud válida de A a B según el código.
Los intercambios B a C y C a A siguen la misma secuencia. En ambos mecanismos
se registran cuatro eventos: envío de la solicitud, recepción de la solicitud,
envío del acuse y recepción del acuse. Las marcas $L_s$, $L_r$ y $L_a$ son los
valores del reloj en esos eventos; los valores reales se verifican a partir de
los CSV en el apartado de Lamport.

\begingroup
\shorthandoff{<>}
\begin{figure}[htbp]
\centering
\begin{tikzpicture}[font=\small,>=Stealth]
\node at (0,0) {Cliente de nodo-A};
\node at (9,0) {Servidor de nodo-B};
\draw[dashed] (0,-0.3) -- (0,-6.6);
\draw[dashed] (9,-0.3) -- (9,-6.6);
\node[anchor=west] at (0.15,-0.75) {Incrementar reloj: $L_s$};
\draw[->] (0,-1.8) -- node[above] {Longitud + JSON: identificador, emisor, $L_s$, texto} (9,-1.8);
\node[anchor=east] at (8.85,-2.6) {Leer trama y validar};
\node[anchor=east] at (8.85,-3.35) {$L_r=\max(L_B,L_s)+1$};
\node[anchor=east] at (8.85,-4.1) {Incrementar para enviar el acuse: $L_a$};
\draw[->] (9,-5.2) -- node[above] {Longitud + JSON: mismo identificador, B, $L_a$, ACK} (0,-5.2);
\node[anchor=west] at (0.15,-6) {Validar acuse y actualizar reloj};
\end{tikzpicture}
\caption{Secuencia TCP/JSON con encuadre definido por la aplicación.}
\label{fig:nixon-secuencia-tcp}
\end{figure}

\begin{figure}[htbp]
\centering
\begin{tikzpicture}[font=\small,>=Stealth]
\node at (0,0) {Cliente de nodo-A};
\node at (4.5,0) {Stub gRPC de A};
\node at (10,0) {Servicio de nodo-B};
\foreach \x in {0,4.5,10} {\draw[dashed] (\x,-0.3) -- (\x,-5.4);}
\draw[->] (0,-0.9) -- node[above] {EnviarMensaje(Mensaje)} (4.5,-0.9);
\draw[->] (4.5,-1.6) -- node[above] {Protobuf sobre HTTP/2 y TCP} (10,-1.6);
\node[anchor=east] at (9.85,-2.2) {Validar y registrar recepción};
\node[anchor=east] at (9.85,-2.75) {Actualizar reloj; formar acuse};
\draw[->] (10,-3.45) -- node[above] {Acuse serializado} (4.5,-3.45);
\draw[->] (4.5,-4.2) -- node[above] {Acuse} (0,-4.2);
\node[anchor=west] at (0.15,-4.9) {Validar identificador, emisor y reloj; registrar recepción};
\end{tikzpicture}
\caption{Secuencia de la llamada unaria gRPC del proyecto.}
\label{fig:nixon-secuencia-grpc}
\end{figure}
\endgroup
